\begin{problem}{even하게 익은 SCON}{standard input}{standard output}{1 second}{256 megabytes}

SCON과 SCCC를 바라보던 문성이는 흥미로운 규칙을 발견했다! 바로 두 문자열에서 문자 `S'와 `C'의 개수를 합하면 항상 짝수가 된다는 것이다. 이를 본 문성이는 문득 알파벳 대문자로 구성된 길이가 $N$인 문자열 중에서 `S'와 `C'의 개수의 합이 짝수인 문자열이 몇 개나 될 지 궁금해졌다.

문성이의 궁금증을 해결해 주자!

\InputFile
첫째줄에 문자열의 길이를 의미하는 자연수 $N(\le 1\,000\,000)$이 주어진다.

\OutputFile
길이 $N$이고 모든 문자가 대문자인 문자열 중 `S', `C'의 개수의 합이 짝수인 문자열의 개수를 $1\,000\,000\,007(=10^9+7)$로 나눈 나머지를 출력하라.

\Examples

\begin{example}
\exmpfile{example.01}{example.01.a}%
\exmpfile{example.02}{example.02.a}%
\end{example}

\end{problem}

